% !TeX root = ../../Book5.tex
% 纲要标题：### 19.1 形式特征标
% 纲要位置：工作记录/计划纲要.md，第 4682 行。
% 纲要细目（保留原定写作范围）：
%
% * 权空间维数
% * 群环中的形式特征标
% * Weyl 群对称性与对偶简单模的最高权
%

\section{形式特征标}
\label{b5:sec:1901}

最高权已经给出了有限维简单模的完整编号，但尚未告诉我们一个模有多大，某个较低权出现多少次，两个简单模的张量积又怎样分解。这些问题可以统一写成权的生成函数。这里的指数符号只记权的加法，不先引入 Lie 群或指数函数。

\subsection{权空间维数}
\label{b5:subsec:1901:01}

本章固定有限维复半单 Lie 代数 \(\mathfrak g\)、Cartan 子代数 \(\mathfrak h\)、正根集 \(\Phi^+\)、权格 \(P\) 和 Weyl 群 \(W\)。内积 \((\, ,\,)\) 始终是 Killing 型在 \(\mathfrak h^*\) 上诱导的型，并限制到根的实空间上使用其正定性。记
\[
\rho=\frac12\sum_{\alpha\in\Phi^+}\alpha.
\]
由\cref{b5:prop:rho-coroot}，\(\rho(h_i)=1\)，所以 \(\rho\in P\)，并处于支配室的内部。我们不另外把长根统一缩放到长度平方二；若读者使用这种归一化，Casimir 特征值和递推两边必须同时作相同缩放。

\begin{definition}\label{b5:def:weight-multiplicity}
设 \(\mathfrak g\) 为有限维复半单 Lie 代数，\(\mathfrak h\) 为其 Cartan 子代数，\(P\subseteq\mathfrak h^*\) 为相应权格。设 \(V\) 为有限维复 \(\mathfrak g\) 模。对 \(\mu\in P\)，以 \(V_\mu\) 表示相对于 \(\mathfrak h\) 的权空间，记
\[
m_V(\mu)=\dim_{\mathbb C}V_\mu,
\]
称为权 \(\mu\) 的\term[quan-chongshu]{权重数}[weight multiplicity]；不是权时规定其为零。另选正根，对相应最高权 \(\lambda\) 的有限维简单模 \(L(\lambda)\)，将其权重数简记为 \(m_\lambda(\mu)\)。
\end{definition}
\symidx{\(m_V(\mu),m_\lambda(\mu)\)：模与简单最高权模的权重数}

最高权分类给出 \(m_\lambda(\lambda)=1\)，而权支集包含于 \(\lambda-Q_+\)。这只是初值与位置限制，不能据此断言其余重数也为一。例如单 Lie 代数的伴随模在零权处的重数是秩；\(\mathfrak{sl}_3\) 的伴随模就有一个二维零权空间。

\subsection{群环中的形式特征标}
\label{b5:subsec:1901:02}

\begin{definition}\label{b5:def:formal-character}
设 \(\mathfrak g\) 为有限维复半单 Lie 代数，\(\mathfrak h\) 为其 Cartan 子代数，\(P\subseteq\mathfrak h^*\) 为相应权格，\(V\) 为有限维复 \(\mathfrak g\) 模。权格群环 \(\mathbb Z[P]\) 以 \(e^\mu\)、\(\mu\in P\) 为自由交换群基，乘法由
\(e^\mu e^\nu=e^{\mu+\nu}\) 给出，其元素只有有限多个非零系数。以 \(m_V(\mu)=\dim_{\mathbb C}V_\mu\) 表示权重数，令
\[
\operatorname{ch}V=\sum_{\mu\in P}m_V(\mu)e^\mu\in\mathbb Z[P].
\]
所得群环元素称为 \(V\) 的\term[xingshi-tezhengbiao]{形式特征标}[formal character]。此处的 \(\operatorname{ch}\) 记录 Cartan 权；它与第七章的 Frobenius 特征映射作用于不同的对象，不能把两种映射直接等同。
\end{definition}
\symidx{\(\mathbb Z[P],e^\mu\)：权格群环及其形式指数基}
\symidx{\(\operatorname{ch}V\)：Lie 代数权模的形式特征标}

\begin{proposition}\label{b5:prop:formal-character-operations}
设 \(\mathfrak g\) 为有限维复半单 Lie 代数，\(\mathfrak h\) 为其 Cartan 子代数，\(P\) 为相应权格，\(V,U\) 为有限维复 \(\mathfrak g\) 模。以 \(m_V(\mu)=\dim_{\mathbb C}V_\mu\) 表示权重数，所有形式特征标均相对于 \(\mathfrak h\) 取值于 \(\mathbb Z[P]\)。则
\[
\operatorname{ch}(V\oplus U)=\operatorname{ch}V+\operatorname{ch}U,
\qquad
\operatorname{ch}(V\otimes U)=\operatorname{ch}V\,\operatorname{ch}U,
\qquad
\operatorname{ch}(V^*)=\sum_\mu m_V(\mu)e^{-\mu}.
\]
短正合列的中项特征标为两端特征标之和。两个有限维模同构，当且仅当形式特征标相同。
\end{proposition}
\begin{proof}
直和按权空间相加。张量作用满足
\(h(v\otimes u)=hv\otimes u+v\otimes hu\)，所以
\((V\otimes U)_\gamma=\bigoplus_{\mu+\nu=\gamma}V_\mu\otimes U_\nu\)，得到乘法。对偶作用为 \((hf)(v)=-f(hv)\)，所以权变号。权空间取值在短正合列上正合：可以使用同时对角化，或用在有限权集上的插值投影逐项取权。

最后，由 Weyl 完全可约性，只须说明简单模的形式特征标线性无关。若有限整数线性组合
\(\sum_\lambda a_\lambda\operatorname{ch}L(\lambda)=0\)，在 \(a_\lambda\ne0\) 的有限集合中选一个对权偏序极大的 \(\lambda\)。其他最高权的特征标不含 \(e^\lambda\)，否则 \(\lambda\le\nu\) 与极大性矛盾。该项系数就是 \(a_\lambda m_\lambda(\lambda)=a_\lambda\)，不可能为零。故简单特征标独立，分解重数由特征标唯一确定。
\end{proof}

\subsection{Weyl 群对称性}
\label{b5:subsec:1901:03}

令 \(W\) 作用于群环，规定 \(w(e^\mu)=e^{w\mu}\)。反射权空间的维数对称性由简单根方向的 \(\mathfrak{sl}_2\) 理论给出。

\begin{proposition}\label{b5:prop:formal-character-weyl-invariant}
设 \(\mathfrak g\) 为有限维复半单 Lie 代数，\(\mathfrak h\) 为其 Cartan 子代数，\(P\) 为相应权格，\(W\) 为相应 Weyl 群。设 \(V\) 为有限维复 \(\mathfrak g\) 模，以 \(m_V(\mu)=\dim_{\mathbb C}V_\mu\) 表示相对于 \(\mathfrak h\) 的权重数，令 \(W\) 在 \(\mathbb Z[P]\) 上按 \(w(e^\mu)=e^{w\mu}\) 作用。则
\[
m_V(w\mu)=m_V(\mu)\quad(\mu\in P,\ w\in W),
\qquad
\operatorname{ch}V\in\mathbb Z[P]^W.
\]
\end{proposition}
\begin{proof}
有限权集与根向量移权公式使每个 \(e_i,f_i\) 幂零：反复加上或减去同一个非零根不能永远留在有限集内。于是\cref{b5:lem:integrable-weight-symmetry}适用，给出每个简单反射保持权重数；它们生成 \(W\)，故结论成立。
\end{proof}

对偶作用把每个权变为其相反数，而最高权还取决于正根的选择。因此，确定对偶模的最高权，需要先找出原模中对同一权偏序最低的权。

\begin{proposition}[对偶模的最高权]\label{b5:prop:dual-highest-weight}
设 \(\mathfrak g\) 为有限维复半单 Lie 代数，固定 Cartan 子代数 \(\mathfrak h\) 及相应根系的一组正根 \(\Phi^+\)。令 \(P^+\) 为支配整权集，\(Q_+\) 为简单根的非负整数组合所成的锥，\(W\) 为 Weyl 群，\(w_0\in W\) 为把全部正根变成负根的\term[zui-chang-yuan]{最长元}[longest element]。对 \(\lambda\in P^+\)，设 \(L(\lambda)\) 为相应有限维简单最高权模。则 \(w_0\lambda\) 是 \(L(\lambda)\) 的最低权，其重数为一，且
\[
L(\lambda)^*\cong L(-w_0\lambda),
\]
其中对偶取左作用 \((xf)(v)=-f(xv)\)。
\end{proposition}
\symidx{\(w_0\)：Weyl 群的最长元}
\begin{proof}
由\cref{b5:thm:weyl-chambers-length}，存在 Weyl 群元素将基本室送到其相反的室，即把全部正根变成负根；长度公式说明它具有最大长度。若 \(u,v\) 都把全部正根变成负根，则 \(u^{-1}v\) 保持全部正根，长度为零，所以 \(u=v\)。而 \(w_0^{-1}\) 也把全部正根变成负根，故 \(w_0^{-1}=w_0\)，即 \(w_0^2=1\)。特别地，\(w_0Q_+=-Q_+\)。

对偶模简单：若 \(F\subseteq L(\lambda)^*\) 是非零真子模，其零化子
\[
F^\perp=\{\,v\in L(\lambda)\mid f(v)=0\;\text{对所有}\;f\in F\,\}
\]
由有限维性也是非零真子空间。对 \(v\in F^\perp\)、\(x\in\mathfrak g\) 及 \(f\in F\)，有 \(f(xv)=-(xf)(v)=0\)，故它还是原模的子模，与原模简单矛盾。

由\cref{b5:prop:highest-weight-support}，原模的权都在 \(\lambda-Q_+\) 中；由\cref{b5:prop:formal-character-weyl-invariant}，权集与重数在 \(W\) 下稳定。对任意权 \(\mu\)，有 \(\lambda-w_0\mu\in Q_+\)。作用 \(w_0\) 后得到 \(w_0\lambda-\mu\in-Q_+\)，所以 \(\mu-w_0\lambda\in Q_+\)。权 \(w_0\lambda\) 确实出现，且重数与最高权 \(\lambda\) 相同，等于一。这证明最低权的结论。

由\cref{b5:prop:formal-character-operations}，对偶模的权为原模各权的负数，故其所有权都不高于 \(-w_0\lambda\)，而该权空间一维。每个正根算子都把它送到不存在的更高权空间，因而将它零化。所以该空间的任一非零向量都是最高权向量；它由对偶模的简单性生成整个模，再由\cref{b5:thm:highest-weight-classification}得到所述同构。
\end{proof}

对 \(\mathfrak{sl}_2\)，令 \(\omega(h)=1\)，则根为 \(2\omega\)。最高权 \(m\omega\) 的特征标为
\[
e^{m\omega}+e^{(m-2)\omega}+\cdots+e^{-m\omega},
\]
在唯一的非平凡反射 \(\omega\mapsto-\omega\) 下不变。高秩时，单个简单反射仍从相应的 \(\mathfrak{sl}_2\) 链得到对称性；需要多个反射一起作用，才控制全部权。
